\chapter{文件 I/O、FITS 与科学数据入口}

\section{本章导读：读入文件之前，先读懂数据}

很多初学者会把数据读取看成一行代码：CSV 用 \texttt{read\_csv}，FITS 用 \texttt{fits.open}，文本用 \texttt{open}。这种理解会让学生很快写出能运行的示例，却很难支撑真实科研。因为程序能打开文件，只说明字节被读进了内存；它并不说明字段含义正确、单位正确、缺失值处理正确、header 可信、数据来源可追踪。

本章的核心观点是：\textbf{I/O 是科学数据工作流的第一道质量控制}。从文件到科学结论，中间至少经过下面几层：

\begin{center}
文件路径 $\rightarrow$ 文件格式 $\rightarrow$ 字段/数组 $\rightarrow$ 类型转换 $\rightarrow$ 单位解释 $\rightarrow$ 数据质量 $\rightarrow$ 科学使用边界
\end{center}

如果前面某一层没有被检查，后面的图表、统计量和模型都可能建立在错误输入上。后续 Part II 会讲星表、图像、光谱、时间序列和实验数据；本章先建立最底层习惯：读文件时不仅要问“能不能读”，还要问“读出来的东西是否能被解释”。

\section{学习目标与训练目标}

学完本章后，你应该能够：

\begin{itemize}
    \item 区分文本、CSV、FITS 等常见科学数据格式的基本用途；
    \item 使用 Python 内置文件接口和 \texttt{csv.DictReader} 读取小型文本和表格；
    \item 解释字段契约、header 契约、metadata、provenance 和 data array 的含义；
    \item 对读入数据做路径、字段、类型、单位、缺失值和样本数检查；
    \item 理解 FITS 为什么适合天文图像和表格；
    \item 写出一份最小数据读入检查记录，说明数据从哪里来、如何读入、哪些值进入后续分析。
\end{itemize}

训练目标是完成一个小型数据入口任务：读入一份星表 CSV 和一个教学版 FITS 图像，做字段/header 检查、缺失值检查、最小统计和背景修正示例，最后交付一份 data intake record。

\section{贯穿本章的微型项目}

本章使用一个小型“星表 + 图像”项目作为情境。它不追求真实巡天复杂度，而是训练学生把数据入口做稳。

\begin{examplebox}[最小数据入口项目]
\begin{lstlisting}[style=terminal]
data_intake_demo/
  data/
    stars_demo.csv
    target_cutout.fits
    README.md
  notebooks/
    inspect_data_inputs.ipynb
  scripts/
    check_data_inputs.py
  results/
    data_intake_record.txt
\end{lstlisting}
\end{examplebox}

\texttt{stars\_demo.csv} 可以包含目标编号、坐标、星等和质量标记；\texttt{target\_cutout.fits} 可以是一张小型教学图像。你要做的不是复杂分析，而是回答：

\begin{itemize}
    \item 文件是否在预期路径；
    \item 表格字段是否完整；
    \item 数值字段是否能转换成数值；
    \item 缺失值和质量标记如何处理；
    \item FITS header 是否包含必要元数据；
    \item 图像数组尺寸和数值范围是否合理；
    \item 这一步检查结果如何写入记录。
\end{itemize}

\section{科学文件不只是字节容器}

任何文件首先都是磁盘上的字节序列。程序需要知道如何解释这些字节：是文本、逗号分隔表格、压缩文件、图像，还是 FITS。科研数据格式的意义就在于给这些字节附加结构；常见格式与入门检查见表~\ref{tab:ch06_file_formats}。

\begin{table}[htbp]
\centering
\begin{tabular}{p{0.18\textwidth}p{0.32\textwidth}p{0.36\textwidth}}
\toprule
格式 & 适合什么 & 入门时必须检查什么 \\
\midrule
文本日志 & 运行记录、错误信息、说明 & 编码、行数、关键词、时间戳 \\
CSV & 小型表格、教学数据、结果摘要 & 字段名、分隔符、缺失值、单位说明 \\
FITS & 天文图像、表格和 header 元数据 & HDU、header、数组形状、单位、WCS 线索 \\
二进制数组 & 大型数值数组或中间缓存 & 形状、dtype、生成程序、版本 \\
\bottomrule
\end{tabular}
\caption{常见科研文件格式及入门阶段必须检查的结构线索。}
\label{tab:ch06_file_formats}
\end{table}

初学阶段最重要的不是掌握所有格式，而是形成共同检查习惯：先确认格式，再解释结构，最后才把数值用于科学判断。

\section{文本 I/O：最小但不能随意}

Python 的内置文件接口适合处理日志、说明文件和小型文本输出。

\begin{examplebox}[写入一个说明文件]
\begin{lstlisting}[style=pythonstyle]
from pathlib import Path

note_path = Path("results/readme_note.txt")
note_path.parent.mkdir(exist_ok=True)
note_path.write_text(
    "target_id: T042\ncomment: suspicious background level\n",
    encoding="utf-8",
)
\end{lstlisting}
\end{examplebox}

这里有三个入门习惯：

\begin{itemize}
    \item 使用 \texttt{Path} 表达路径；
    \item 显式创建输出目录；
    \item 写文本时说明编码。
\end{itemize}

读取文本也应先抽样：

\begin{examplebox}[读取文本前几行]
\begin{lstlisting}[style=pythonstyle]
log_path = Path("logs/run.log")
if log_path.exists():
    lines = log_path.read_text(encoding="utf-8").splitlines()
    print(lines[:5])
\end{lstlisting}
\end{examplebox}

文本文件常常承担“解释”和“证据”角色。它们看似不如 FITS 或数组高级，但在复现工作流中非常重要。

\section{CSV：透明但脆弱的教学入口}

CSV 是教学中最常用的数据格式之一，因为它透明、容易查看、容易生成。但 CSV 也很脆弱：它通常不强制单位，不强制字段类型，不强制缺失值规则，也不天然保存复杂元数据。

\begin{examplebox}[读取一个最小 CSV]
\begin{lstlisting}[style=pythonstyle]
import csv
from pathlib import Path

data_path = Path("data/stars_demo.csv")
with data_path.open(newline="", encoding="utf-8") as handle:
    rows = list(csv.DictReader(handle))

print(len(rows))
print(rows[0].keys())
\end{lstlisting}
\end{examplebox}

读入后不应立刻画图或建模。第一步应该是数据入口检查：

\begin{enumerate}
    \item 行数是否接近预期；
    \item 字段是否完整；
    \item 关键字段是否有单位说明；
    \item 数值字段能否转换；
    \item 缺失值如何表示；
    \item 质量标记是否需要筛选。
\end{enumerate}

\section{字段契约：表格必须说清自己的结构}

一个表格能被读入，不代表它被正确理解。科研表格至少需要一份字段契约。

\begin{defbox}[字段契约]
字段契约是对表格结构的最小约定，包括必需字段、字段含义、单位、类型转换规则、允许的缺失值表示和质量标记含义。
\end{defbox}

例如一份星表教学数据可以按表~\ref{tab:ch06_field_contract_basic} 规定：

\begin{table}[htbp]
\centering
\begin{tabular}{p{0.24\textwidth}p{0.22\textwidth}p{0.36\textwidth}}
\toprule
字段 & 单位或类型 & 含义 \\
\midrule
\texttt{source\_id} & string & 目标编号 \\
\texttt{ra\_deg} & degree & 赤经 \\
\texttt{dec\_deg} & degree & 赤纬 \\
\texttt{mag} & mag & 观测星等 \\
\texttt{mag\_err} & mag & 星等不确定度 \\
\texttt{quality} & string & 教学质量标记 \\
\bottomrule
\end{tabular}
\caption{星表教学数据的最小字段契约示例。}
\label{tab:ch06_field_contract_basic}
\end{table}

还可以把同一份契约写成“含义加检查”的形式，见表~\ref{tab:ch06_field_contract_checks}：

\begin{table}[htbp]
\centering
\begin{tabular}{p{0.20\textwidth}p{0.24\textwidth}p{0.18\textwidth}p{0.26\textwidth}}
\toprule
字段 & 含义 & 单位 & 最小检查 \\
\midrule
\texttt{source\_id} & 对象编号 & 无 & 是否唯一、是否缺失 \\
\texttt{ra\_deg} & 赤经 & degree & 是否在 \([0,360)\) \\
\texttt{dec\_deg} & 赤纬 & degree & 是否在 \([-90,90]\) \\
\texttt{mag} & 星等 & mag & 是否可转为数值 \\
\texttt{quality} & 质量标记 & 无 & 是否属于允许集合 \\
\bottomrule
\end{tabular}
\caption{字段契约中字段含义、单位和最小检查的对应关系。}
\label{tab:ch06_field_contract_checks}
\end{table}

这张表比“数据有这些列”更进一步：它说明每一列怎样被解释、应该怎样检查、什么情况下不能进入后续分析。字段契约写不清时，后面的模型和图表都会缺少地基。

字段检查可以很小：

\begin{examplebox}[检查必需字段]
\begin{lstlisting}[style=pythonstyle]
required = {"source_id", "ra_deg", "dec_deg", "mag", "mag_err", "quality"}
available = set(rows[0].keys()) if rows else set()
missing = required - available

if missing:
    raise ValueError(f"Missing required fields: {sorted(missing)}")
\end{lstlisting}
\end{examplebox}

这一步越早做越好。字段名错误如果拖到绘图或模型训练阶段才暴露，通常更难定位。

\section{类型转换：字符串还不是科学量}

CSV 读入后，每个值通常都是字符串。字符串 \texttt{"15.3"} 看起来像数值，但还不能直接参与数值计算；字符串 \texttt{""}、\texttt{"NA"}、\texttt{"nan"} 也不一定会自动变成你想要的缺失值。

\begin{examplebox}[最小数值解析函数]
\begin{lstlisting}[style=pythonstyle]
def parse_float(value, field_name):
    text = value.strip()
    if text in {"", "NA", "nan", "None"}:
        return None
    try:
        return float(text)
    except ValueError as exc:
        raise ValueError(f"{field_name} is not numeric: {value!r}") from exc
\end{lstlisting}
\end{examplebox}

有了这类函数，数据入口就从“读到了字符串”推进到“知道哪些值可以作为数值使用”。

\begin{warnbox}[不要把缺失值自动填成零]
零是一个有物理含义的数值，缺失值表示未知、不适用或未观测。把缺失值随手填成 0，可能会制造虚假的亮源、零误差或不存在的物理边界。
\end{warnbox}

\section{单位：数值必须连接到物理含义}

科学数据中的数字只有和单位、坐标系、时间系统或仪器定义连接起来，才有物理含义。数值 \texttt{1.0} 可能是 1 degree、1 radian、1 arcsec、1 count、1 Jy，也可能只是归一化单位。

入门阶段至少要养成两条习惯：

\begin{itemize}
    \item 字段名中尽量包含单位线索，例如 \texttt{ra\_deg}、\texttt{flux\_jy}；
    \item 在数据 README 或 intake record 中写明单位，不要只靠代码变量名。
\end{itemize}

单位检查不一定要很复杂。很多错误可以通过数量级发现：

\begin{examplebox}[最小范围检查]
\begin{lstlisting}[style=pythonstyle]
ra_values = [parse_float(row["ra_deg"], "ra_deg") for row in rows]
bad_ra = [value for value in ra_values if value is not None and not (0 <= value < 360)]

if bad_ra:
    raise ValueError(f"RA outside [0, 360): {bad_ra[:3]}")
\end{lstlisting}
\end{examplebox}

\section{缺失值与质量标记：数据状态也是数据}

缺失值不一定表示错误。它可能来自未观测、低信噪比、上游筛选、交叉匹配失败或教学数据故意设置。质量标记也是数据的一部分，它说明某些行是否适合进入后续分析。

\begin{examplebox}[按质量标记筛选]
\begin{lstlisting}[style=pythonstyle]
clean_rows = []
for row in rows:
    mag = parse_float(row["mag"], "mag")
    mag_err = parse_float(row["mag_err"], "mag_err")
    if row["quality"] == "good" and mag is not None and mag_err is not None:
        clean_rows.append({**row, "mag": mag, "mag_err": mag_err})

print(len(rows), len(clean_rows))
\end{lstlisting}
\end{examplebox}

筛选不是纯技术动作，而是科学假设。你必须记录：哪些行被排除，为什么排除，排除后样本是否仍然代表原问题。

\section{原始数据与清洗数据：不要覆盖入口证据}

数据清洗的目标不是把原始数据“修成正确”，而是从原始数据中生成一份有明确规则的可分析数据。这个区别很重要。原始数据承担证据角色，清洗数据承担分析入口角色；两者应该能互相追溯，但不应该互相覆盖。

一个安全的最小流程是：

\begin{enumerate}
    \item 原始文件放在 \texttt{data/raw/} 或课程指定位置，默认只读；
    \item 清洗脚本读取原始文件，按照字段契约和质量规则筛选；
    \item 清洗结果写入 \texttt{data/processed/} 或 \texttt{results/}；
    \item 同时写出排除记录，说明哪些行因为什么原因没有进入分析；
    \item 在 data intake record 中写明原始输入、清洗输出和脚本入口。
\end{enumerate}

如果你直接在原始 CSV 中删除坏行，短期看文件更干净，长期却丢失了判断依据。别人无法知道坏行原来是什么、为什么被删、是否删错。更可靠的做法是让清洗规则存在于脚本和记录中，让原始数据保持可追溯。

\begin{warnbox}[清洗数据不是静默改文件]
任何会改变样本集合的操作都应留下记录。尤其是删除缺失值、剔除异常值、按质量标记筛选、单位转换和重命名字段，都应该能在脚本或 data intake record 中找到理由。
\end{warnbox}

\section{从文件到可分析数据的最小算法}

把 CSV 数据入口压缩成算法，可以写成：

\begin{enumerate}
    \item 确认路径存在；
    \item 读取文件并检查行数；
    \item 检查必需字段；
    \item 根据字段契约做类型转换；
    \item 检查单位和数值范围；
    \item 处理缺失值和质量标记；
    \item 保存清洗后样本数、排除理由和清洗输出路径；
    \item 写入 data intake record，并保留原始输入的追溯线索。
\end{enumerate}

这条算法会在后续章节不断重复。无论是星表、FITS 图像、光谱、时间序列还是机器学习特征表，入口检查永远不应该省略。

\section{逐行验证：不要让坏行悄悄混进去}

字段存在不代表每一行都可用。更可靠的做法是把验证拆到每一行：

\begin{examplebox}[逐行验证骨架]
\begin{lstlisting}[style=pythonstyle]
def validate_row(row):
    problems = []

    try:
        ra = float(row["ra_deg"])
        dec = float(row["dec_deg"])
    except ValueError:
        problems.append("coordinate_not_numeric")
        return problems

    if not (0.0 <= ra < 360.0):
        problems.append("ra_out_of_range")
    if not (-90.0 <= dec <= 90.0):
        problems.append("dec_out_of_range")
    if row["quality"] not in {"good", "suspect", "bad"}:
        problems.append("unknown_quality_flag")

    return problems
\end{lstlisting}
\end{examplebox}

逐行验证的目的不是让代码变长，而是让排除规则可解释。正式记录中不应只写“清洗后剩余 83 行”，还应说明哪些行因为什么原因被排除。

\section{Provenance：数据从哪里来}

即使字段和数值都正确，你仍然要知道数据从哪里来。这个信息称为 provenance。

\begin{defbox}[Provenance]
Provenance 指数据来源、版本、生成脚本、处理步骤和关键假设。它回答“这份数据为什么可信，别人如何追溯它”。
\end{defbox}

一个最小 provenance 记录可以包含：

\begin{examplebox}[最小 provenance 字典]
\begin{lstlisting}[style=pythonstyle]
provenance = {
    "source_file": str(data_path),
    "format": "csv",
    "required_fields": sorted(required),
    "row_count": len(rows),
    "clean_row_count": len(clean_rows),
    "quality_rule": "quality == good and mag/mag_err present",
}
\end{lstlisting}
\end{examplebox}

在课程项目中，provenance 不一定需要复杂数据库。先做到 README、manifest、脚本名、commit hash 和检查记录，就已经能显著提高可复查性。

\section{FITS：天文学为什么需要专门格式}

FITS 是天文学长期使用的数据格式。它常用于保存图像、表格和元数据。一个 FITS 文件通常由一个或多个 HDU 组成；每个 HDU 可以包含 header 和 data。

\begin{defbox}[FITS Header]
FITS header 是一组关键字和值，记录观测时间、仪器、曝光时间、单位、坐标线索、数据维度等元数据。它告诉你 data array 应该如何被解释。
\end{defbox}

FITS 的重要性在于它把数组和元数据放在同一个文件结构中。对于天文图像，仅有像素矩阵是不够的；你还需要知道这些像素来自哪个目标、什么曝光、什么单位、怎样映射到天空坐标。

\section{HDU：一个 FITS 文件里可能不止一份数据}

HDU 是 Header/Data Unit 的缩写。一个 FITS 文件可以包含多个 HDU，例如表~\ref{tab:ch06_fits_hdu_checks} 中的几类：

\begin{table}[htbp]
\centering
\begin{tabular}{p{0.18\textwidth}p{0.30\textwidth}p{0.36\textwidth}}
\toprule
HDU & 可能内容 & 需要检查什么 \\
\midrule
Primary HDU & 主图像或空数据加主 header & 图像形状、曝光时间、目标名 \\
Image extension & 误差图、掩膜图、权重图 & 是否与主图像同形状 \\
Table extension & 源表、检测结果、质量表 & 字段名、单位、行数 \\
\bottomrule
\end{tabular}
\caption{FITS 文件中常见 HDU 类型及读取时的检查重点。}
\label{tab:ch06_fits_hdu_checks}
\end{table}

因此，\texttt{hdul[0].data} 并不总是你真正需要的科学数据。读取 FITS 时先运行 \texttt{hdul.info()}，就是为了确认文件内部结构，而不是盲目取第一个数组。

\section{最小 FITS 读取流程}

使用 \texttt{astropy.io.fits} 可以读取 FITS：

\begin{examplebox}[读取 FITS 的最小流程]
\begin{lstlisting}[style=pythonstyle]
from astropy.io import fits
from pathlib import Path

fits_path = Path("data/target_cutout.fits")
with fits.open(fits_path) as hdul:
    hdul.info()
    header = hdul[0].header
    image = hdul[0].data

print(image.shape)
print(header.get("EXPTIME"))
\end{lstlisting}
\end{examplebox}

入门阶段至少检查：

\begin{itemize}
    \item 有多少个 HDU；
    \item 图像数组形状是什么；
    \item 数组中是否存在 NaN；
    \item header 中是否有曝光时间、目标名、单位或坐标线索；
    \item 数据值范围是否明显异常。
\end{itemize}

\section{Header 契约：FITS 也需要入口检查}

CSV 有字段契约，FITS 有 header 契约。它们本质上是同一个思想：数据本体必须和解释它的说明一起检查。

\begin{examplebox}[检查关键 header]
\begin{lstlisting}[style=pythonstyle]
required_cards = ["NAXIS", "NAXIS1", "NAXIS2"]
missing_cards = [card for card in required_cards if card not in header]

if missing_cards:
    raise ValueError(f"Missing FITS header cards: {missing_cards}")
\end{lstlisting}
\end{examplebox}

不同数据产品需要检查的 header 不同。教学图像可以只检查数组维度；真实天文图像还可能需要检查 \texttt{EXPTIME}、\texttt{FILTER}、\texttt{BUNIT}、WCS 关键字和数据质量扩展。

\section{WCS：像素坐标到天空坐标的桥}

天文图像中的一个像素位置，例如 \((x,y)=(120,80)\)，本身只是数组索引。科学解释往往还需要知道这个像素对应天空上的哪个方向。WCS（World Coordinate System）就是描述像素坐标和天空坐标之间关系的元数据系统。

入门阶段不需要推导完整 WCS 变换，但要知道最小检查包括：

\begin{itemize}
    \item header 中是否存在 \texttt{CTYPE}、\texttt{CRVAL}、\texttt{CRPIX}、\texttt{CD} 或 \texttt{CDELT} 等线索；
    \item 图像是否已经被裁剪或重采样；
    \item 画图时坐标轴是像素坐标还是天空坐标；
    \item 后续测光或匹配是否需要真正的天球坐标。
\end{itemize}

如果没有 WCS，图像仍然可以用于教学数组练习，但不能随意声称某个像素对应某个天球位置。这就是数据入口检查和科学解释边界之间的关系。

\section{图像数组：矩阵不是自动等于图像}

FITS 图像读入后通常是二维数组。数组中的每个元素可能表示 count、count rate、flux、背景修正后的值或归一化强度。你不能只看形状，还要理解数值含义。

\begin{examplebox}[最小数组检查]
\begin{lstlisting}[style=pythonstyle]
import math

height, width = image.shape
flat_values = [float(value) for row in image for value in row]
finite_values = [value for value in flat_values if math.isfinite(value)]

print(height, width)
print(min(finite_values), max(finite_values))
\end{lstlisting}
\end{examplebox}

真实项目中通常会使用 NumPy 更高效地完成这些检查。本章示例保留朴素写法，是为了让学生看清“数组检查到底在检查什么”。

\section{背景修正光度的最小公式}

为了说明“读入图像后仍然需要解释数值”，考虑一个最小孔径测光思想。假设孔径内有 \(N_{\rm ap}\) 个像素，像素值为 \(I_i\)，背景环带估计出的每像素背景为 \(b\)，则净信号可以写成：

\[
S_{\rm net} = \sum_{i \in \mathrm{aperture}} I_i - N_{\rm ap} b.
\]

这个公式很简单，但它说明数据入口和科学解释之间的连接：

\begin{itemize}
    \item 图像数组中的像素值是什么单位；
    \item 哪些像素进入孔径；
    \item 背景 \(b\) 如何估计；
    \item 是否存在坏像元或饱和像元；
    \item 结果是否应该除以曝光时间。
\end{itemize}

因此，FITS 读取不是图像分析的结束，而只是图像分析的入口。

\section{曝光时间与单位：同一个 count 不一定同义}

两个图像中同样的像素值 \texttt{1000}，如果曝光时间不同，物理含义就可能不同。一个最小的 count rate 可以写成：

\[
R = \frac{C}{t_{\rm exp}},
\]

其中 \(C\) 是计数，\(t_{\rm exp}\) 是曝光时间。若 header 中没有可靠的 \texttt{EXPTIME}，你就不能轻易比较不同图像的亮度。若 \texttt{BUNIT} 表示数据已经是 count rate，再除一次曝光时间反而会出错。

这说明 header 不是附属说明，而是公式的一部分。你使用什么公式，取决于数据值已经表达了什么单位。

\section{把入口检查写成记录}

本章最重要的交付物不是一段漂亮代码，而是一份数据读入检查记录。

\begin{examplebox}[最小数据读入检查记录]
\begin{lstlisting}[style=terminal]
Dataset: data_intake_demo
CSV file: data/stars_demo.csv
FITS file: data/target_cutout.fits
Reader notebook: notebooks/inspect_data_inputs.ipynb
Reader script: scripts/check_data_inputs.py
Commit:

CSV checks:
- row count:
- required fields:
- missing fields:
- clean row count:
- unit notes:

FITS checks:
- HDU count:
- image shape:
- required header cards:
- finite pixel fraction:
- value range:

Decisions:
- rows used:
- rows excluded:
- reason for exclusion:

Limits:
- teaching data only
- WCS not fully validated
\end{lstlisting}
\end{examplebox}

这份记录会进入后续 evidence cards 中的 Data 和 Reproducibility 部分。它让别人知道你不是“随手读了一个文件”，而是做了入口质量控制。

\section{配套 notebook 如何训练}

本章配套 notebook 会把正文内容落到一个可运行流程中：

\begin{enumerate}
    \item 创建或读取小型 CSV；
    \item 检查字段契约；
    \item 做类型转换和缺失值统计；
    \item 按质量标记筛选；
    \item 读取教学 FITS 文件；
    \item 检查 header 和数组形状；
    \item 计算最小数组统计量；
    \item 生成 data intake record。
\end{enumerate}

训练重点不是让你记住某个库函数，而是让你建立数据入口习惯。以后遇到新数据集，你应该先问“契约是什么”，再问“怎么分析”。

\section{常见错误与诊断}

\begin{warnbox}[数据 I/O 中的典型问题]
\begin{itemize}
    \item 文件路径正确，但字段名和预期不一致；
    \item 读入后立刻画图或建模，没有检查行数、字段和缺失值；
    \item 把字符串当作数值使用；
    \item 把缺失值填成 0；
    \item 不记录单位，导致度、弧度、角秒或像素混用；
    \item 只看 FITS 数组，不看 header；
    \item 数据来源、生成脚本和版本没有记录；
    \item 清洗掉多少样本没有写入报告。
\end{itemize}
\end{warnbox}

诊断数据入口问题时，先不要改模型。先回到路径、字段、类型、单位、缺失值和来源这六件事。很多模型异常，其实是数据入口阶段埋下的问题。

\section{AI 助手如何帮助你}

AI 助手可以帮助你设计字段契约、阅读 header、生成检查代码、整理 data README 和总结异常值。但它不能替你确认数据的科学含义。尤其是单位、质量标记、坏像元和筛选条件，必须回到数据文档和科学问题。

一个好的提示词是：

\begin{lstlisting}[style=terminal]
请根据这份 CSV 的字段名和前 5 行样本，帮我设计一个字段契约。
请指出必需字段、可能的单位、缺失值风险和最小检查代码。
\end{lstlisting}

对于 FITS，可以这样问：

\begin{lstlisting}[style=terminal]
请根据下面的 FITS header 关键词，判断哪些信息足够解释图像数组，
哪些信息仍然缺失。请不要替我假设单位或坐标系统。
\end{lstlisting}

\section{练习与交付}

\begin{exercisebox}[基础练习]
读取一份小型 CSV，打印行数和字段名。写出字段契约，至少包含字段含义、单位和是否允许缺失。
\end{exercisebox}

\begin{exercisebox}[类型练习]
为两个数值字段写类型转换函数，统计无法转换或缺失的行数。解释为什么这些行是否应该进入后续分析。
\end{exercisebox}

\begin{exercisebox}[FITS 练习]
读取一个教学 FITS 文件，列出 HDU 信息、图像形状、关键 header 卡片和像素值范围。说明这张图像还能不能直接用于物理解释。
\end{exercisebox}

\begin{exercisebox}[交付练习]
写一份 \texttt{results/data\_intake\_record.txt}，记录 CSV 和 FITS 的路径、字段/header 检查、缺失值、筛选决策、输出文件和已知限制。
\end{exercisebox}

\section{本章小结}

数据读取不是把文件塞进程序，而是建立数据和科学解释之间的第一座桥。一个可靠的数据入口流程必须检查路径、格式、字段、类型、单位、缺失值、header、数组形状和来源记录。只要这一步做扎实，后续绘图、统计、机器学习和 capstone 项目才有可信输入。忽略 I/O 质量控制，后面的模型越复杂，错误越难发现。
